\section{Introduction}

In collaborative interaction, the most informative action is often not an answer but a question.

A tutor locates the boundary of a learner's understanding. A research assistant discovers accepted
assumptions. A decision-support agent identifies the uncertainty blocking progress.

In each case, assistance depends on an internal estimate of the user and an action that improves it.

This is more demanding than describing a mental state after its evidence has been supplied.

The agent must decide which uncertainty matters, acquire evidence under a budget, revise its estimate,
and let that estimate change its next action.

User modeling is therefore part of the control loop, not an annotation attached to a completed dialogue.

Existing evaluations illuminate parts of this problem but do not identify the full loop.

Machine Theory-of-Mind (ToM) benchmarks infer beliefs, intentions, or emotions from fixed narratives
and conversations \citep{le2019revisiting,gandhi2023bigtom,kim2023fantom,xu2024opentom,chen2024tombench}.

Their controlled questions make inference measurable, yet the benchmark chooses the evidence. The model
is not accountable for an uninformative question or unresolved uncertainty.

Interactive social-agent environments restore action, but usually score goal completion, persuasion, or
task success \citep{zhou2024sotopia,mou2025agentsense,hwang2026infusing,wang2026tomelp,zhou2026tomswe}.

These outcomes combine user modeling with planning, domain competence, language generation, and
strategy. Success may reflect an accurate user model, a generic policy, or environmental shortcuts.

The inference--application gap in explicit ToM evaluation makes this ambiguity consequential
\citep{gu2024simpletom}.

The resulting gap is one of \emph{measurement}, not simply coverage.

Static tasks remove evidence acquisition. Interactive tasks restore action but often hide the latent
user state from the evaluator.

Neither setting tests whether an agent improves an explicit user model by choosing what evidence to
request.

We study the following question:

\begin{quote}
\emph{Can an agent actively infer a user's hidden knowledge state through limited natural-language interaction, and use that evolving estimate to choose subsequent diagnostic actions?}
\end{quote}

We call the target behavior \emph{functional ToM for knowledge-grounded interaction}.

\emph{Functional} means that a user model earns credit through an infer--update--act loop, not through a
claim that the model understands the user.

\emph{Knowledge-grounded} narrows the latent state to reviewed, diagnosable concepts. The definition
identifies useful behavior without asserting a human-like cognitive mechanism.

A valid test requires three commitments.

First, the latent state must be explicit and directly scoreable. Second, evidence must be selected under
a finite budget. Third, the interaction must be auditable.

Without these commitments, task success, exhaustive questions, leakage, or simulator artifacts can
masquerade as diagnosis.

\knowact turns these commitments into a controlled benchmark.

A public graph defines the shared hypothesis space, while a hidden user map defines the target. The
agent asks diagnostic questions under a fixed turn budget and submits a prediction for every node.

The \sagesim simulator grounds a question before reading hidden state, abstracts only licensed
epistemic content, and realizes it separately. The final map is scored deterministically.

This design makes competing explanations testable.

Fixed and random policies test whether interaction alone is sufficient. Working-map and graph ablations
test whether explicit state guides action.

Simulator variations test proxy dependence. Evidence references separate dialogue-supported
reconstruction from accurate guesses.

It also creates a technical agent problem that the benchmark alone does not solve.

The agent must represent node uncertainty, ground open answers in rubrics, use graph relations without
copying mastery, value the next probe, and decide when evidence is not worth its cost.

We propose MapProbe as an explicit evidence--belief--probe loop. It exposes evidence, beliefs, target
plans, and stopping decisions so a final score can be attributed to a mechanism.

Our contributions are therefore design and protocol contributions rather than empirical claims at this stage:

\begin{itemize}[leftmargin=1.5em]
    \item We formulate active knowledge-state diagnosis as an identifiable target for functional user modeling, with \sagesim as a scoped and independently testable observation model.
    \item We specify MapProbe, a graph-aware, evidence-backed reconstruction agent whose belief update, diagnostic target policy, question realization, and stopping rule are independently inspectable and ablatable.
    \item We predefine a paired evaluation that separates active acquisition, explicit state use, graph use, simulator robustness, repeated reliability, and human transfer instead of collapsing them into one leaderboard score.
\end{itemize}

The visibility boundary, full-map shell, and simple prompted agent baseline are implemented.

An experimental ECDA slice implements node-marginal belief updates and multi-candidate question
utility. The full MapProbe design and comparative experiments remain incomplete or unrun.

The paper therefore does not provide evidence that an agent exhibits the target behavior. This
distinction is deliberate: \knowact makes the claim falsifiable before it is made.
